\section{Threats to validity}
\label{sec:threats}

\textbf{Internal validity.}
Scheduling variants are compared at byte-identical predictions, so
accuracy is not a confound. Timing confounds we controlled are run
order (rotated), thermal state (10\,s pauses; junction temperature
logged; on the Jetson, maximum 67.4\,$^\circ$C in application runs and 69.2\,$^\circ$C in the \texttt{trtexec} sweep, far from throttling --- the Rubik CPU runs, by contrast, sit at the 95\,$^\circ$C thermal limit and show a quantified 17\% drift, \cref{sec:rubik}), page-cache state
(warm reads before each phase) and clock state (stored and restored,
verified after each phase). Remaining confounds are the small number of
repeats per cell and the 20\,Hz sampling of power and clocks, which
cannot resolve sub-50\,ms clock changes (\cref{sec:diagnosis}). The
memory (EMC) clock has its own governor. A direct devfreq lock attempt
did not hold on this release --- a silent-failure mode documented
independently on the same board family~\cite{kang2026emc} --- so in
the fixed-frequency sweep, where only the GPU domain is pinned, we
logged the memory clock and bounded its influence to 2\% of the fitted
slope (\cref{sec:diagnosis}). \texttt{jetson\_clocks} does hold it
(verified by readback, \cref{sec:setup}), so the locked-versus-default
comparisons include any memory-clock effect. Jetson campaign runs last 9--203\,s, and the
cross-board accuracy sweeps run longer (the Rubik CPU sweep takes
${\approx}26$\,min); we did not test
deployment-length steady state, and the ramp visible in
\cref{fig:traces} means the first seconds of a run differ from the rest. The
Ultralytics baseline uses its own engine export. Its AP agrees with
ours to $3\times10^{-5}$, but engine tactics may differ slightly.

\textbf{Construct validity.}
\texttt{VDD\_IN} is the module's main input rail, not whole-system power: it excludes carrier-board loads such as the fan, USB peripherals and storage regulators. It still overstates the energy of the SoC compute rails alone, and the idle cost of locking clocks includes the higher memory clock as well as the CPU and GPU clocks. We did not validate the on-board sensor against an external power meter; our comparisons rely on its consistency between runs rather than on its absolute accuracy. Our ``latency'' starts when a JPEG
file is available to the pipeline; decode is part of the measured work. Camera capture, image-signal-processor (ISP)
processing and output
transmission are excluded. Process CPU time is reported per image and
includes all threads of the process.

\textbf{External validity.}
The main study uses one Jetson Orin Nano Super with one software
stack; \cref{sec:crossplatform} partially offsets this with a
replication on a Hailo-8 NPU board and a Qualcomm board measured on both its CPU and its Hexagon NPU, but
other Jetson modules, power modes, JetPack releases, hardware JPEG decode~\cite{chen2026nvjpeg} or DeepStream pipelines may shift the magnitudes, and other vendors' governors may behave differently. On the Pi~5 the Hailo-8 sits on the unmonitored 5\,V rail, so Pi energies cover the SoC only, and its HEF's on-chip NMS uses a 0.2 score threshold, so Pi accuracy is not comparable with the other boards. On the Rubik Pi~3 there is no usable power sensor; the NPU stack runs through a third-party bundled runtime (\texttt{qai-appbuilder}) rather than the vendor SDK, and its primary INT8 context ends in the standard head with host-side NMS (a fused NMS-free head was measured as a configuration comparison and costs 7 AP on this toolchain, \cref{sec:rubikhtp}), so its accuracy is not comparable with the float stacks; the QNN userspace wedged under concurrent inference calls, so our NPU pipeline serializes accelerator calls and overlaps host stages only; and multi-threaded CPU inference emits rows in a nondeterministic order, so prediction equivalence on the CPU stack is shown on canonicalized (sorted) rows --- identical in all 24 ladder cells --- rather than by a byte-level hash. Our pipelines are written in Python. A native C++ or DeepStream pipeline would shorten the host stages and could reach the GPU-bound regime with fewer decode workers. We also used batch size 1 throughout, as a camera application would. Batch size and GPU frequency are known to interact~\cite{nabavinejad2022batching}; we did not measure batching. The compression screening used a short recovery schedule
on a development split with one seed and \texttt{opt0} TensorRT
builds. It is evidence about that
budget, not a general verdict on pruning or quantization.

\section{Conclusion}
\label{sec:conclusion}

For NMS-free detectors on an embedded GPU, the deployed cost of a
model is not a property of the model alone. Under the stock governor,
the same YOLO26n engine costs 4.2\,ms or 13.2\,ms depending on how busy
the application keeps the GPU. A simple duty-cycle model, with
inverse-clock engine time, a load-band governor and an idle-plus-marginal
energy split, describes the measurements. Restructuring only the host pipeline, without locking clocks and at unchanged COCO accuracy, raised throughput $2.4\times$ over the best sequential pipeline with locked clocks and $6.9\times$ over the standard predictor, and cut energy per image $2.6\times$ relative to the latter. The CPU governor takes part in the same loop: a sleeping CUDA wait lowers the host's clock and, through the host, the GPU's. A per-task utilization clamp, set without root privileges, recovers most of this governor-mediated loss; on the full two-worker schedule the wait policy itself costs nothing, and only the synchronization boundary (24\% for stream synchronization) remains, at locked and default clocks alike. At camera rates
below the pipeline's capacity the picture changes again: idle power
then sets energy per frame, and the schedule mainly sets latency ---
until the rate approaches capacity, where the schedule decides energy
once more (117--123\,mJ per processed frame for the pipelines that
keep up at 60\,fps, 151--194\,mJ for those that do not). The right schedule therefore
depends on the arrival rate. The same restructuring changed
the value of every model-level decision we examined. Lower resolution
gave only $2$--$4$\% throughput in the optimized pipeline ($6$--$10$\%
with the garbage collector disabled), against
$27$\% in the naive one. Compression and routing
gave less than scheduling. The pipeline effect was larger than the step from YOLO26n to the $1.6\times$ slower YOLO26s. On a Raspberry Pi~5 with a Hailo-8 NPU the same host restructuring gives $2.7\times$ ($3.3\times$ with the runtime's wait policy replaced) with identical predictions, even though the accelerator exposes no frequency control, while on a Qualcomm board the ladder flattens to $1.3\times$ on the CPU --- where thermal management becomes the binding constraint --- and on its Hexagon NPU the CPU governor becomes the largest term, because the scheduler down-clocks the core that feeds the accelerator. We hope that these results and the
accompanying data encourage edge-detection research to report
end-to-end, governor-aware measurements alongside engine benchmarks.
